\section{NVSHMEM}
NVSHMEM e uma biblioteca de comunicao para GPUs NVIDIA baseada no modelo OpenSHMEM. 
O objetivo eh oferecer um espaco de memoria global entre GPUs, onde cada elemento de 
processamento possue uma memoria local, que pode ser enderecado globalmente atraves de
operacoes de comunicacao. O principal diferencial do NVSHMEM e o motivo de ser utilizado
nesse trabalho eh que as operacoes podem ser emitidas diretamente de um kernel CUDA,
sem envolver a CPU \cite{hsunvshmem}.

As unidades de processamento sao divididasaaa em PEs (Process Unit). A biblioteca define
regioes de memoria simetricas, essas regioes sao alocadas coletivamente por todas as PEs.
A simetria permite que uma PE, calcule o endereco logico global correspondente a um
endereco logico local, e um indentificador de PE. \leal{\sout{O NVSHMEM fornece operacoes remotas
como leituras e escritas, e atomicos.} explicar essa parte ou tirar, oq é uma operação atomica}

Pelas suas caracteristicas a biblioteca NVSHMEM e relevante para algoritmos irregulares.
Em aplicacoes densas ou estruturadas, a comunicacao pode ser prevista e agrupada
em mensagens grandes. Em grafos, por outro lado, os destinos das atualizacoes
sao descobertos durante a execucao. Ao expandir uma fronteira de BFS, uma thread
pode visitar uma aresta, identificar que o vertice vizinho pertence a outro PE e
precisar atualizar o estado desse vertice remoto. Com NVSHMEM, essa atualizacao
pode ser expressa diretamente no kernel, usando uma operacao remota ou atomica,
sem retornar primeiro ao host para empacotar mensagens \cite{hsunvshmem}.

\section{MPI}
MPI (\textit{Message Passing Interface}) e um padrao de comunicacao para
programas paralelos em memoria distribuida. Em um programa MPI, a execucao e
organizada em processos independentes, chamados normalmente de ranks. Cada rank
possui seu proprio espaco de enderecamento e troca dados com os demais por meio
de chamadas explicitas de comunicacao. Esse modelo e adequado para clusters,
onde a memoria de uma maquina ou processo nao pode ser acessada diretamente por
outro processo sem uma operacao de comunicacao definida \cite{mpistandard}.

\leal{ \sout{
O modelo de MPI se baseia em troca de mensagens. Em comunicacoes ponto a ponto,
um rank envia dados e outro rank recebe esses dados. Operacoes como
\texttt{MPI\_Send} e \texttt{MPI\_Recv} representam esse padrao basico. Tambem
existem versoes nao bloqueantes, como \texttt{MPI\_Isend} e
\texttt{MPI\_Irecv}, que permitem iniciar uma transferencia e continuar a
execucao enquanto a comunicacao progride. Esse tipo de comunicacao e importante
em aplicacoes paralelas porque pode permitir sobreposicao entre computacao e
transferencia de dados.

MPI tambem define operacoes coletivas, nas quais todos os ranks de um
comunicador participam. Exemplos comuns incluem \texttt{MPI\_Barrier}, usada
para sincronizacao, \texttt{MPI\_Bcast}, usada para difundir dados de um rank
para os demais, e \texttt{MPI\_Reduce} ou \texttt{MPI\_Allreduce}, usadas para
combinar valores produzidos localmente por cada rank. Em algoritmos iterativos,
essas operacoes aparecem frequentemente para calcular condicoes globais de
parada, agregar estatisticas ou garantir que todos os processos avancem para a
proxima fase de execucao.
}}